\documentclass[12pt]{article}

\usepackage{comment}
\begin{document}

\centerline{\textbf{Gurkeerat Singh}}


\centerline{Homework 01, MORALLY Due Feb 10}
\newcommand{\implies}{\Rightarrow}

(The last question IS honors HW 01 and graded separatel but handed in
with this HW.) 


\begin{enumerate}

\item
(0 points)
When is the midterm (give date and time)?
If you cannot take the midterm on that day then tell 
Professor Gasarch by Feb 10. 
\begin{itemize}
    \item \textbf{The midterm is on March 31st, 2025. Since it's in class it'll go from 9:30 to 10:45 A.M.}
\end{itemize}

\bigskip

\newpage

\item
(40 points)
\begin{enumerate}
\item 
(20 points)
Give a Propositional Formula on four variables that has exactly
three satisfying assignments.
Give the satisfying assignments.
\textbf {A formula on 4 variables that has exactly 3 satisfying assignments is $ (a \wedge b \wedge (c \vee d)) $. The 3 satisfying assignements would look like this:}
\begin{itemize}
    \item A = \textbf{TRUE}, B = \textbf{TRUE}, C = \textbf{TRUE}, D = \textbf{TRUE}
    \item A = \textbf{TRUE}, B = \textbf{TRUE}, C = \textbf{TRUE}, D = \textbf{FALSE}
    \item A = \textbf{TRUE}, B = \textbf{TRUE}, C = \textbf{FALSE}, D = \textbf{TRUE}
\end{itemize}


\bigskip

\item 
(20 points)
Give a Propositional Formula on four variables that has exactly
13 satisfying assignments.
ADVICE: You could do it in a messy way BUT there is an EASIER way.
Both are full credit, but try to find the easier way. 

\textbf{Since there are 4 variables, there are a total of 16 input combinations; therefore, an easier way to find this propositional formula, rather than finding 13 satisfying assignments, would be to find 3 assignments that don't satisfy the formula. We already have a formula with 3 satisfying assignments; therefore negating the whole formula would give us 13 satisfying assignments. So $ \neg (a \wedge b \wedge (c \vee d)) $ has 13 satisfying assignments.
}

\bigskip

\end{enumerate}

\newpage

\item
(30 points)
Recall that a {\it 2CNF} formula is of the form 

$$C_1 \wedge \cdots \wedge C_k$$

where each $C_i$ is an $\vee$ of two literals. 

Recall that  {\it 2SAT} is the following problem: Given a 2CNF formula, determine
if it is satisfiable. 

Recall that I said in class that 2SAT is in P.

\begin{enumerate}
\item
(0 points but you need to do this for Part 2)
Look on the web or ChatGPT or whatever you want for an ALGORITHM that
solves 2SAT in polynomial time.
\item
(30 points) Describe the algorithm in English so that someone who
can't read code can follow it. (You do not need to show that its in polynomial time.) 
\begin{itemize}
    \item \textbf{When we deal with 2SAT, we are dealing with a formula made up of some number of clauses, with each clause having the form $ (x \vee y) $. These clauses are logically equivalent to $ (\neg x \rightarrow y) \wedge (\neg y \rightarrow x)$, giving us two different implications. Since disjunction is commutative, not x would imply y and not y would imply x, meaning that they both work for the clause form. Once this is understood we can create a graph with nodes, giving each variable x two nodes, one x and one $\neg x$. From there we can connect each $\neg a$ to a b, and each $\neg b$ to an a, showing how $\neg a$ implies b and $\neg a$ implies b.
    Once we are done connecting nodes, look at the most strongly-connected parts - areas where every node can reach every other node in the group. If x and $\neg x$ are in the same area, we know that the formula is unsatisfiable, as both x and $\neg x$ can't be true. If there are none of these scenarios, the formula is SAT.}
\end{itemize}
\item
(0 points and just for your own benefit and part of a long term extra credit problem) 
Code up the 2SAT algorithm. 
This is used in every SAT algorithm. I may have an extra credit project in this course
where you code up a SAT solver. 
\end{enumerate}


\newpage

\item
(30 points)
Show that, for all $n\ge 1$, there exists a formula on $n$ variables that is satisfied by
exactly $n$ satisfying assignments. Give the satisfying assignments.
(This is NOT by induction. Just give
the Formula. You may use DOT DOT DOT (that is ``$\dots$") though it should be clear what you mean.)

\begin{itemize}
    \item \textbf{If we have a function with n variables, a formula that made exactly one variables true and the rest false would have n different satisfying assignments, one variable true for each n. The formula for this would look like this: } \vspace{0.3cm} $$ f(x_1, \dots , x_n) = \bigvee_{i=1}^n \left ( x_i \wedge \bigwedge_{j \ne i} \neg x_j \right)$$ \vspace{0.3cm} 
    \textbf{This formula ensures that only one x is true at a time. If we have n different x's, there are n different satisfying arguments for this problem.}
    
\end{itemize}

\newpage
\item 
{\bf Honors HW 1} This is graded separately as Honors HW01.

There are 3 people in a room.

Two are truth tellers.

One is normal (can tell the truth or lie).

They all know about each other. 

Show how by asking them YES-NO questions you can determine which one is the normal. 

\textbf{Let's define 3 people, two truth-tellers and one normal person as persons A, B, and C. We can start by asking person B if person A is normal. If person B says that person A IS normal, either A or B is normal, since B can be a truth-teller and A can be normal, or B can be normal and lie, saying that A is normal. Therefore, we know C is a truth teller and we can ask him if A is normal. His answer (yes leading to B and no leading to A) will help us determine who is normal. In the case person B says that person A IS NOT normal, person B may either be a truth-teller or a normal once again. If person B is a truth-teller, that suggests person C is the normal. Otherwise, person B is the normal, and he's the person we're looking for. So when person B says person A is NOT normal, we turn to person A(a guaranteed truth-teller) to ask if person B is normal, where the answer (yes leads to B and no leads to C) will help us understand who is normal and who isn't.}

\bigskip

\end{enumerate}
\end{document}
